% Generated draft: requires mathematical and source review.
\documentclass[11pt]{article}
\usepackage[margin=25mm]{geometry}
\usepackage{fontspec}
\setmainfont{Latin Modern Roman}
\usepackage{amsmath,amssymb,mathtools,unicode-math}
\setmathfont{Latin Modern Math}
\usepackage{enumitem,needspace,xurl,hyperref}
\setlength{\parindent}{0pt}
\setlength{\parskip}{5pt}
\setlength{\emergencystretch}{3em}
\newcommand{\N}{\mathbb{N}}
\newcommand{\Z}{\mathbb{Z}}
\newcommand{\Q}{\mathbb{Q}}
\newcommand{\R}{\mathbb{R}}
\newcommand{\C}{\mathbb{C}}
\newcommand{\F}{\mathbb{F}}
\newcommand{\A}{\mathbb{A}}
\newcommand{\PP}{\mathbb{P}}
\newcommand{\E}{\mathbb{E}}
\newcommand{\eps}{\varepsilon}
\newcommand{\dd}{\,\mathrm{d}}
\begin{document}
\section*{UnsolvedMath ID: 30002817}
\textbf{Source problem code:} catalogue entry
\subsection{Definitions and mathematical statement}
Let \(D_1,D_2\subset\R^2\) be smooth convex plane domains, labelled so that \(D_1\subseteq D_2\). An oriented line is an affine line equipped with a unit direction vector. Write
\[
\mathcal L=\{\ell:\ell\text{ is an oriented line meeting both }D_1\text{ and }D_2\}.
\]
For a line directed by \(v\) that has a finite transverse forward exit point \(x\in\partial D_i\), the usual billiard map \(T_i\) sends that line to the oriented line through \(x\) with direction
\[
v'=v-2\langle v,n_i(x)\rangle n_i(x),
\]
where \(n_i(x)\) is the outward unit normal. This defines the reflection convention wherever such an exit exists.

The supplied question asks whether, for every pair of nested smooth convex plane domains in the intended class, commutation of their billiard maps on the intended common space of oriented lines implies that their boundaries are confocal ellipses:
\[
T_1\circ T_2=T_2\circ T_1
\quad\Longrightarrow\quad
\exists f_1,f_2\in\R^2\ \exists a_1,a_2>\tfrac12\|f_1-f_2\|\ 
\forall i\in\{1,2\},\qquad
\partial D_i=\{x\in\R^2:\|x-f_1\|+\|x-f_2\|=2a_i\}.
\]
The displayed conclusion expresses a common pair of foci; allowing \(f_1=f_2\) includes concentric circles as the coincident-focus case.

\emph{Unresolved formulation qualifications.} The source does not specify whether the domains must be bounded, distinct, strictly convex, or strictly nested, nor does it specify the precise smoothness class or treatment of tangencies. For the displayed composition equality, \(T_1(T_2(\ell))\) requires \(T_2(\ell)\) to lie in the domain of \(T_1\), and conversely. The source does not establish that \(\mathcal L\) is invariant under either map or specify whether equality means equality of partial maps, equality wherever both compositions exist, or equality on a designated invariant phase space. These interpretations are not silently identified here. The source also does not specify whether coincident foci are permitted. Consequently, the displayed implication records the intended question but requires these qualifications to be resolved before it is a fully determined conjecture.
\subsection{Short English statement}
The question asks whether two nested smooth convex planar billiard tables whose reflection maps commute must have elliptical boundaries with the same foci. The admissible tables and the domain of the composition identity need clarification. The supplied dated curation labels the problem partially solved and reports a result for piecewise-analytic domains; these claims have not been independently verified.
\subsection{Sources}
\begin{itemize}
\item Supplied local UnsolvedMath registry record, permanent ID \(30002817\), title \emph{Commuting Billiard Maps and Confocal Ellipses}, exact source problem code \texttt{catalogue entry}, in the Oberwolfach Reports Open Problems collection. The record gives proposed year 2015 and names no individual proposer.
\item Supplied report provenance: \emph{Discrete Differential Geometry} (2015), Oberwolfach Reports, pp.~661--729, DOI \texttt{10.4171/owr/2015/13}. \url{https://doi.org/10.4171/owr/2015/13}. The supplied background describes Oberwolfach Reports as documenting mathematical workshops and their open-problem discussions; no problem-specific page is supplied.
\item Supplied literature source labelled ``Glutsyuk, billiard ball maps'': \url{https://arxiv.org/abs/1405.5990}.
\item Supplied literature source labelled ``OWR Report 13/2015'': \url{https://doi.org/10.4171/owr/2015/13}.
\item Supplied literature source labelled ``Recent billiard-map results'': \url{https://arxiv.org/abs/2607.23853}.
\item Literature assessment dated exactly \texttt{2026-08-21}: the supplied curation labels the problem \texttt{partially\_solved}, states that the confocal-ellipse conclusion is proved for piecewise-analytic nested billiard domains, and states that the unrestricted smooth or general extracted formulation remains unresolved. This is supplied dated curation derived from the OWR record, not an independently verified literature assessment; the cited claims and time-sensitive status require verification.
\item The registry labels the statement \texttt{corrected\_verified} and reports manual semantic, typographic, and LaTeX review of a meaning-preserving cleanup. This supplied editorial label does not resolve the mathematical formulation issues identified here or constitute independent mathematical review.
\end{itemize}
\end{document}
